\section{Scelte Tecnologiche e Scelte Implementative}

\subsection{Ambiente di Sviluppo e Strumenti Utilizzati}
Lo sviluppo del progetto è stato condotto su sistema operativo Linux. Per la gestione del runtime e l'isolamento delle dipendenze è stato impiegato l'interprete Python~3.14 all'interno di un ambiente virtuale dedicato (\texttt{venv}), garantendo la piena riproducibilità delle librerie specificate nel file \texttt{requirements.txt}.

Il tracciamento del codice sorgente e la collaborazione sono stati gestiti tramite il sistema di controllo di versione Git, mentre la relazione è redatta in \LaTeX{} per la formattazione di equazioni, algoritmi e tabelle.

\subsection{Linguaggio di Programmazione e Librerie}
Python è stato adottato come linguaggio di riferimento per la chiarezza sintattica del codice e per la disponibilità di ottime librerie dedicate all'analisi delle prestazioni.
Nello specifico, l'implementazione si avvale sia di moduli della libreria standard sia di pacchetti esterni:
\begin{itemize}
	\item \textbf{Libreria standard di Python:}
	      \begin{itemize}
		      \item \texttt{dataclasses} e \texttt{typing}: impiegati per formalizzare le entità cardine del dominio (\texttt{Player}, \texttt{Role}, \texttt{Scenario}, \texttt{RosterSolution}, \texttt{BenchmarkMetric}) secondo un modello tipizzato e coerente. L'attributo \texttt{frozen=True} garantisce l'immutabilità delle istanze, prevenendo alterazioni accidentali dei dati durante l'esecuzione degli algoritmi.
		      \item \texttt{pathlib} e \texttt{json}: per gestire i percorsi dei file in modo indipendente dal sistema operativo e per caricare gli scenari di test in formato JSON.
		      \item \texttt{time.perf\_counter}: per misurare i tempi di calcolo con la massima precisione offerta dal sistema.
		      \item \texttt{tracemalloc}: per misurare la memoria allocata esclusivamente dalle strutture dati dell'algoritmo (come le matrici dinamiche), escludendo l'overhead dell'interprete Python e la memoria occupata dai moduli importati.
	      \end{itemize}
	\item \textbf{Librerie esterne:}
	      \begin{itemize}
		      \item \texttt{pandas} e \texttt{numpy}: impiegate nell'infrastruttura di benchmark per l'aggregazione, l'ordinamento e l'esportazione dei risultati su file CSV.
		      \item \texttt{matplotlib}: impiegata per la generazione dei grafici di scalabilità temporale e spaziale e della tabella comparativa riassuntiva.
	      \end{itemize}
\end{itemize}

\subsection{Architettura del Software e Modularità}
Per mantenere i file e il codice ordinati e separare chiaramente la logica risolutiva dal caricamento dei dati e dalle misurazioni, il progetto all'interno di \texttt{src/} è suddiviso in tre cartelle principali, ciascuna con uno scopo specifico:
\begin{enumerate}
	\item \textbf{\texttt{core/}}: racchiude la definizione degli oggetti di dominio (\texttt{models.py}) e la logica di caricamento e validazione degli scenari da file JSON (\texttt{loader.py}).
	\item \textbf{\texttt{algorithms/}}: ospita l'implementazione algoritmica pura delle funzioni e delle relative varianti. Le procedure operano direttamente su matrici e liste di dati, rimanendo completamente disaccoppiate da operazioni di I/O o di cronometraggio.
	\item \textbf{\texttt{benchmark/}}: dedicata all'analisi comparativa tra gli algoritmi. I file \texttt{profiler.py}, \texttt{exporter.py} e \texttt{plotter.py} gestiscono rispettivamente la rilevazione di tempi e memoria, il salvataggio dei risultati in CSV e la produzione dei grafici.
\end{enumerate}
L'esecuzione complessiva è gestita dallo script \texttt{main.py}, che esegue uno dopo l'altro gli scenari, raccoglie i risultati e produce i file di output.

\subsection{Scelte Implementative Rilevanti}
Durante la trasposizione implementativa delle formulazioni teoriche descritte nel Compito~1, sono state introdotte alcune scelte mirate all'efficienza computazionale e all'accuratezza numerica:

\subsubsection{Discretizzazione del Budget e dei Costi}
Poiché i compensi contrattuali e il budget sono espressi in unità monetarie reali nell'ordine dei milioni di dollari, è stata definita una costante di scala $\text{COST\_UNIT} = 100\,000$. Ciò consente di mappare il budget $X$ e i costi dei singoli atleti su valori interi discreti:
\begin{equation}
	B = \lfloor X / \text{COST\_UNIT} \rfloor, \qquad c_{i, k} = \lfloor \mathit{cost}_{i, k} / \text{COST\_UNIT} \rfloor
\end{equation}
Tale discretizzazione permette l'accesso alle celle della tabella dinamica tramite indici interi diretti in tempo $\Theta(1)$, prevenendo al contempo errori di approssimazione dovuti all'aritmetica floating-point nel controllo dei vincoli di bilancio.

\subsubsection{Rappresentazione delle Tabelle Dinamiche}
Per la memorizzazione degli stati della programmazione dinamica è stato preferito l'uso di liste bidimensionali native Python preallocate (\texttt{list} di \texttt{list}) rispetto a strutture esterne come gli array NumPy. Data la presenza di accessi ed aggiornamenti scalari frequenti all'interno di cicli iterativi annidati, le liste native di CPython risultano più efficienti, non introducendo l'overhead di chiamata dell'interfaccia C-API. Inoltre, la propagazione dello stato base ($m[i, b] \gets m[i-1, b]$) è demandata allo slicing nativo \texttt{m[i][:] = m[i-1][:]}, eseguito internamente in C a basso livello.

\subsubsection{Trade-off nella Ricostruzione della Rosa: con e senza Matrice \texorpdfstring{$s$}{s}}
Per la determinazione della rosa ottimale sono state implementate due strategie complementari, delineando un chiaro compromesso tra occupazione spaziale e complessità temporale di ricostruzione:
\begin{itemize}
	\item \textbf{Ricostruzione con matrice $s$:} prevede l'allocazione di una matrice ausiliaria $s[i, b]$ per memorizzare l'indice della scelta ottima a ogni passo. La ricostruzione della rosa si compie risalendo la matrice a ritroso in tempo $\Theta(N)$, con un costo addizionale di $(N+1) \times (B+1)$ interi in memoria primaria.
	\item \textbf{Ricostruzione senza matrice $s$:} evita del tutto l'allocazione della tabella delle decisioni, determinando la soluzione a ritroso direttamente sulla matrice dei rendimenti $m$. Per ciascun ruolo $i$, viene identificato il candidato $p$ che soddisfa la transizione di Bellman $m[i-1, b - c_{i, p}] + v_{i, p} = m[i, b]$ entro una tolleranza numerica $\varepsilon = 10^{-9}$ (necessaria per gestire la rappresentazione floating-point dei WAR). Il tempo di risalita diviene $O(N \cdot P)$, ma la memoria richiesta nella fase di calcolo risulta dimezzata.
\end{itemize}

\subsubsection{Scansione a Ritroso nell'Ottimizzazione a Vettore Singolo}
Nella variante a memoria compressa, la tabella bidimensionale $m$ è ridotta a un singolo vettore di dimensione $B+1$, abbattendo la complessità spaziale da $\Theta(N \cdot B)$ a $\Theta(B)$. Per rispettare il vincolo di mutua esclusività (al più un ingaggio per ruolo), l'aggiornamento degli stati sul budget viene eseguito in senso decrescente da $B$ verso $0$. L'inversione della scansione impedisce che l'elaborazione di un candidato sovrascriva i valori calcolati allo stadio decisionale precedente, garantendo che nessun atleta del medesimo ruolo venga selezionato più volte nel medesimo passaggio.

\subsubsection{Filtraggio Preliminare di Pareto}
Prima del calcolo dinamico, i candidati di ciascun ruolo vengono sottoposti a un filtro di dominanza paretiana. Ordinando gli atleti per costo crescente, vengono conservati unicamente coloro che presentano un WAR strettamente superiore al massimo valore registrato per i candidati di costo inferiore. I profili dominati (aventi costo maggiore o uguale ma WAR minore o uguale) non possono matematicamente migliorare la funzione obiettivo e vengono scartati a priori. Tale potatura riduce la cardinalità effettiva da $P$ a $P'$, accelerando il calcolo nelle varianti che operano direttamente sui valori senza richiedere indici statici in $s$.

\subsection{Metodologia di Benchmarking e Profiling}
Per assicurare rigore metodologico e ripetibilità nell'analisi empirica:
\begin{itemize}
	\item \textbf{Misurazione del tempo di calcolo:} ciascuna esecuzione è ripetuta per $K = 15$ iterazioni, assumendo come metrica il valore minimo. Tale approccio attenua l'influenza del rumore di sistema (context-switch del sistema operativo, interrupt hardware o interventi del garbage collector), approssimando con accuratezza il reale tempo di calcolo computazionale.
	\item \textbf{Misurazione della memoria RAM:} il consumo di memoria è quantificato campionando il picco di byte tracciati tramite \texttt{tracemalloc} tra l'inizio e la fine dell'esecuzione, isolando l'impronta generata dalle matrici e dai vettori dinamici dall'occupazione statica dell'interprete Python.
\end{itemize}
